\article{表示的张量积分解}{Decomposition of Representations into Tensor Products}{D. Flath}{1}{179--183}
\origsection{1}{有限直积的表示}
本文讨论以下经典定理的若干推广：两个有限群的直积的不可约表示（irreducible representation），由两个因子各自的不可约表示确定。我们首先将这一结果推广到更一般的群。
\begin{originaltheorem}{定理 1（Theorem 1）}
设 $G_1,G_2$ 是局部紧全不连通群（locally compact totally disconnected groups），$G=G_1\times G_2$。

（1.1）若 $\pi_i$ 是 $G_i$ 的容许不可约表示（admissible irreducible representation），$i=1,2$，则 $\pi_1\otimes\pi_2$ 是 $G$ 的容许不可约表示。

（1.2）若 $\pi$ 是 $G$ 的容许不可约表示，则存在 $G_i$ 的容许不可约表示 $\pi_i$，使得 $\pi\simeq\pi_1\otimes\pi_2$。各 $\pi_i$ 的等价类由 $\pi$ 的等价类确定。
\end{originaltheorem}
先回顾所需的记号。对于局部紧全不连通群 $G$，其 Hecke 代数（Hecke algebra）$H(G)$ 是 $G$ 上具有紧支集的局部常值复函数组成的卷积代数（convolution algebra）。对 $G$ 的紧开子群（compact open subgroup）$K$，令
\[
e_K=(\operatorname{meas}K)^{-1}\operatorname{ch}_K,
\]
其中 $\operatorname{ch}_K$ 是 $K$ 的示性函数（characteristic function），$\operatorname{meas}$ 是定义 $H(G)$ 的卷积时所用的 Haar 测度。于是 $e_K$ 是 $H(G)$ 的幂等元（idempotent）。子代数 $e_KH(G)e_K$ 记为 $H(G,K)$。光滑 $G$ 模（smooth $G$-module）$W$ 自然成为 $H(G)$ 模；对每个紧开子群 $K\subset G$，空间 $W^K=e_KW$ 是 $H(G,K)$ 模。

证明定理 1 之前，先给出\textbf{不可约性判别准则（Irreducibility Criterion）}：光滑 $G$ 模 $W$ 不可约，当且仅当对 $G$ 的所有紧开子群 $K$，$W^K$ 都是不可约的 $H(G,K)$ 模。\footnote{\textbf{译者注（本次校订）：}准确表述应为：非零光滑模 $W$ 不可约，当且仅当每个 $W^K$ 都是零模或不可约的 $H(G,K)$ 模。不可约表示可以没有 $K$ 不变向量，因而必须允许 $W^K=0$；下面的推论也因此写作 $\dim W^K\leq1$。后文应用这一判别准则时均采用此表述。}
\begin{proof}
若 $U$ 是 $W^K$ 的 $H(G,K)$ 子模，则 $(H(G)\cdot U)^K=U$；由此即可推出判别准则。
\end{proof}
应用这一判别准则时，只需对构成 $G$ 的单位元处邻域基的一族 $K$，检验 $W^K$ 是不可约 $H(G,K)$ 模。
\begin{originaltheorem}{推论（Corollary）}
设 $K$ 是 $G$ 的紧开子群，$H(G,K)$ 交换，$W$ 是容许不可约 $G$ 模。则 $\dim W^K\leq1$。
\end{originaltheorem}
\footnotetext{AMS（MOS）学科分类（1970）：主要分类 22E55。}
\begin{proof}[定理 1 的证明]
对 $G_i$ 的紧开子群 $K_i$ 及光滑 $G_i$ 模 $W_i$（$i=1,2$），有以下自然同构：
\begin{align*}
\text{（i）}\quad&H(G_1\times G_2)\simeq H(G_1)\otimes H(G_2),\\
\text{（ii）}\quad&H(G_1\times G_2,K_1\times K_2)\simeq H(G_1,K_1)\otimes H(G_2,K_2),\\
\text{（iii）}\quad&(W_1\otimes W_2)^{K_1\times K_2}\simeq W_1^{K_1}\otimes W_2^{K_2}.
\end{align*}
结合（iii）与不可约性判别准则，即得（1.1）。

反之，设 $W$ 是容许不可约 $G$ 模。取 $K=K_1\times K_2$，其中 $K_i$ 是 $G_i$ 的紧开子群，$i=1,2$，并使 $W^K\ne0$。空间 $W^K$ 有限维，故由 [2] 第 94 页的推论，存在不可约的 $H(G_i,K_i)$ 模 $W_i^{K_i}$，以及从 $W^K$ 到 $W_1^{K_1}\otimes W_2^{K_2}$ 的 $H(G,K)$ 同构 $a_K$。同样的论证适用于 $K$ 的每个形如 $K'=K'_1\times K'_2$ 的开子群。存在 $H(G_i,K_i)$ 映射
\[
b_i=b_i(K,K'):W_i^{K_i}\longrightarrow W_i^{K_i'},
\]
使下面的图交换：
\[
\begin{tikzcd}
W^K \arrow[r,"a_K"] \arrow[d,hook,"\text{包含}"'] & W_1^{K_1}\otimes W_2^{K_2} \arrow[d,"b_1\otimes b_2"]\\
W^{K'} \arrow[r,"a_{K'}"'] & W_1^{K_1'}\otimes W_2^{K_2'} .
\end{tikzcd}
\]
此外，可以为每对这种类型的紧开子群 $K,K'$ 选择映射 $b_i(K,K')$，使它们构成归纳系（inductive system）。于是 $W\simeq W_1\otimes W_2$，其中
\[
W_i=\varinjlim_{K_i}W_i^{K_i},
\]
而 $W_i$ 是 $G_i$ 的容许不可约表示，$i=1,2$。

$W_i$ 的等价类由 $W$ 的等价类确定，因为 $W$ 在 $G_i$ 上的限制是 $W_i$ 等型的（$W_i$-isotypic）。
\end{proof}
定理 1 的证明只用到了 $G$ 与 $G_i$ 的 Hecke 代数。这启发我们引入\textbf{带幂等元的代数（idempotented algebra）} $(A,E)$：它由代数 $A$ 及幂等元的一个有向族 $E$ 构成，满足 $A=\bigcup_{e\in E}eAe$。$(A,E)$ 的\textbf{容许模（admissible module）}是一个 $A$ 模 $W$，它非退化（nondegenerate），即 $AW=W$，且对所有 $e\in E$，$\dim eW$ 有限。两个带幂等元的代数的张量积自然也带有幂等元。略作调整，定理 1 的证明便可推广到这种情形，给出两个带幂等元代数的张量积之容许不可约模的分类。

\origsection{2}{受限张量积}
阿代尔群（adelic groups）是群的无限受限积（restricted product）；研究其表示，需要用到 [4] 所引入的向量空间的受限张量积（restricted tensor product）。

设 $\{W_v\mid v\in V\}$ 是一族向量空间，$V_0$ 是 $V$ 的有限子集。对每个 $v\in V\setminus V_0$，选取非零向量 $x_v\in W_v$。对包含 $V_0$ 的每个有限子集 $S\subset V$，令 $W_S=\bigotimes_{v\in S}W_v$；若 $S\subset S'$，则定义 $f_S^{S'}:W_S\to W_{S'}$ 为
\[
\bigotimes_{v\in S}w_v\longmapsto
\bigotimes_{v\in S}w_v\otimes\bigotimes_{v\in S'\setminus S}x_v.
\]
$W_v$ 关于 $x_v$ 的受限张量积 $W=\bigotimes_{x_v}W_v$ 定义为 $W=\varinjlim_S W_S$。空间 $W$ 由可写成 $w=\bigotimes w_v$ 的元素张成，其中对几乎所有 $v\in V$，$w_v=x_v$。

有限张量积的常用构造可以自然推广到受限张量积：

（1）给定线性映射 $B_v:W_v\to\widetilde W_v$，且对几乎所有 $v\in V$ 有 $B_vx_v=\widetilde x_v$，则可定义 $B=\bigotimes B_v:W\to\widetilde W$，令 $B(\bigotimes w_v)=\bigotimes B_vw_v$。

（2）给定一族代数 $\{A_v\mid v\in V\}$，且对几乎所有 $v$ 给定非零幂等元 $e_v\in A_v$，则 $A=\bigotimes_{e_v}A_v$ 自然成为一个代数。

（3）若每个 $v\in V$ 的 $W_v$ 都是 $A_v$ 模，且对几乎所有 $v$ 有 $e_v\cdot x_v=x_v$，则 $\bigotimes_{x_v}W_v$ 是 $A$ 模。$W$ 的同构类依赖于 $\{x_v\}$。不过，若 $\{x'_v\}$ 是另一族非零向量，而且对几乎所有 $v$，$x_v$ 与 $x'_v$ 张成 $W_v$ 中的同一条直线，则 $A$ 模 $\bigotimes_{x_v}W_v$ 与 $\bigotimes_{x'_v}W_v$ 同构。

\paragraph*{例 1（Example 1）}无穷多个变量的多项式环 $\C[X_1,X_2,\ldots]$ 同构于 $\bigotimes_{e_i}\C[X_i]$，其中 $e_i$ 是 $\C[X_i]$ 的单位元。

\paragraph*{例 2（Example 2）}设 $G=\prod'_{K_v}G_v$ 是局部紧全不连通群 $G_v$ 关于紧开子群 $K_v$ 的受限积。则 $G$ 本身局部紧且全不连通，并且 $H(G)$ 同构于 $\bigotimes_{e_{K_v}}H(G_v)$。

对每个 $v\in V$，设 $W_v$ 是容许 $G_v$ 模。假设对几乎所有 $v$，$\dim W_v^{K_v}=1$。对几乎所有 $v$ 选取非零向量 $x_v\in W_v^{K_v}$，便可构造 $G$ 模 $W=\bigotimes_{x_v}W_v$。$W$ 的同构类实际上与 $x_v\in W_v^{K_v}$ 的选择无关，称这个表示为各 $W_v$ 的张量积。容易看出，$W$ 容许，且 $W$ 不可约当且仅当每个 $W_v$ 都不可约。$G$ 的容许不可约表示若同构于这样构造的表示，就称为可分解的（factorizable）。
\begin{originaltheorem}{定理 2（Theorem 2）}
假设对几乎所有 $v$，$H(G_v,K_v)$ 交换。则 $G$ 的每个容许不可约表示 $W$ 均可分解，$W\simeq\bigotimes W_v$。各因子 $W_v$ 的同构类由 $W$ 的同构类确定。对几乎所有 $v$，$\dim W_v^{K_v}=1$。
\end{originaltheorem}
\begin{proof}
先取 $G$ 的紧开子群 $K'=\prod K'_v$，分解有限维空间 $W^{K'}$，再沿用定理 1 的论证。
\end{proof}

\origsection{3}{阿代尔群的表示}
设 $G$ 是整体域（global field）$F$ 上的连通约化代数群（connected reductive algebraic group）。令 $A$ 为 $F$ 的阿代尔环（adele ring），$V$ 为 $F$ 的素位（places）集合。阿代尔群 $G(A)$ 同构于受限积 $\prod'_{K_v}G(F_v)$；其中，对所有有限素位 $v$ 都定义了子群 $K_v$，它们是 $G(F_v)$ 的某些极大紧子群。对几乎所有有限素位 $v\in V$，$G$ 在 $F_v$ 上拟分裂（quasi-split group），$K_v$ 是特殊极大紧子群（special maximal compact subgroup）。对这些素位 $v$，$H(G(F_v),K_v)$ 交换，见 [5]。因此，下述定理在函数域情形是定理 2 的特例；数域情形则还需说明“容许表示”的含义。
\begin{originaltheorem}{定理 3（Theorem 3）}
$G(A)$ 的每个容许不可约表示均可分解。各因子在等价意义下唯一。
\end{originaltheorem}
设 $F$ 是数域（number field）。在这种情形，需要给出 $G(A)$ 的容许表示的定义。对每个阿基米德素位（archimedean place）$v\in V$，令 $K_v$ 为 $G(F_v)$ 的极大紧子群，$\mathfrak g_v$ 为 $G(F_v)$ 的实 Lie 代数。令
\[
K_\infty=\prod_{\text{阿基米德 }v}K_v,\qquad
K=\prod_{\text{所有 }v}K_v,\qquad
G_\infty=\prod_{\text{阿基米德 }v}G(F_v).
\]
令 $\mathfrak g_\infty$ 为 $G_\infty$ 的实 Lie 代数，$A_f$ 为 $F$ 的有限阿代尔环。

\paragraph*{定义（Definition）}容许 $G(A)$ 模 $W$ 是同时具有 $(\mathfrak g_\infty,K_\infty)$ 模结构和光滑 $G(A_f)$ 模结构的向量空间 $W$，且满足：

（1）$G(A_f)$ 的作用与 $\mathfrak g_\infty$ 和 $K_\infty$ 的作用交换；

（2）对 $K$ 的每个连续不可约表示同构类 $\gamma$，$W$ 的 $\gamma$ 等型分量（isotypic component）有限维。

定理 3 中，阿基米德素位 $v$ 处的因子应当是容许 $(\mathfrak g_v,K_v)$ 模。只需为每个阿基米德素位 $v$ 构造一个带幂等元的代数，使其容许模对应于容许 $(\mathfrak g_v,K_v)$ 模，便可沿用函数域情形的证明。

设 $G$ 是 Lie 群，$K$ 是紧子群，$\mathfrak g$ 和 $\mathfrak t$ 分别是 $G$ 和 $K$ 的实 Lie 代数。令 $U(\mathfrak g_\C)$ 与 $U(\mathfrak t_\C)$ 为复化 Lie 代数的泛包络代数（universal enveloping algebras）。将 $(\mathfrak g,K)$ 的 Hecke 代数 $H(\mathfrak g,K)$ 定义为 $G$ 上支集包含于 $K$ 的左右 $K$ 有限分布（left and right $K$-finite distributions）组成的代数。它包含代数 $A_K$，后者由 $K$ 上的 $K$ 有限测度构成，视为 $G$ 上的分布。映射 $(X,\mu)\mapsto X*\mu$，取值于 $G$ 上的分布空间，并诱导出向量空间同构
\[
U(\mathfrak g_\C)\otimes_{U(\mathfrak t_\C)}A_K\simeq H(\mathfrak g,K).
\]
以 $A_K$ 的中心幂等元集合为 $E$，$H(\mathfrak g,K)$ 是带幂等元的代数。

设 $(\pi,W)$ 是 $(\mathfrak g,K)$ 模。由公式
\[
(X\otimes\mu)\cdot w=\pi(X)\pi(\mu)w,
\qquad X\in U(\mathfrak g_\C),\ \mu\in A_K,\ w\in W,
\]
空间 $W$ 成为非退化 $H(\mathfrak g,K)$ 模。容易验证，这一构造给出容许 $(\mathfrak g,K)$ 模范畴与容许 $(H(\mathfrak g,K),E)$ 模范畴之间的同构。

\origsection{4}{Hilbert 空间中的分解}
实际应用还需要在 Hilbert 空间框架下建立相应的分解理论。

设 $\{H_v\mid v\in V\}$ 是一族 Hilbert 空间。对几乎所有 $v\in V$，设 $x_v$ 是 $H_v$ 中的单位向量。Hilbert 受限积（Hilbert restricted product）$H=\widehat\bigotimes_{x_v}H_v$ 是一个 Hilbert 空间，可用以下标准正交基来描述。对每个 $v\in V$，令 $P_v$ 为 $H_v$ 的标准正交基；对几乎所有 $v$，要求它包含 $\{x_v\}$。所有符号 $\widehat\bigotimes h_v$ 构成 $H$ 的标准正交基，其中对所有 $v$ 有 $h_v\in P_v$，且对几乎所有 $v$ 有 $h_v=x_v$。前面关于代数受限张量积的各项构造，在 Hilbert 空间中都有相应的版本。

下述定理见 [3] 及 [1]。
\begin{originaltheorem}{定理 4（Theorem 4）}
设 $G(A)$ 如定理 3，$\pi$ 是 $G(A)$ 的连续不可约酉 Hilbert 空间表示（continuous irreducible unitary Hilbert space representation）。则：

（1）存在 $G(F_v)$ 的连续不可约酉 Hilbert 空间表示 $\pi_v$，且几乎所有 $\pi_v$ 均非分歧（unramified），使得 $\pi\simeq\widehat\bigotimes\pi_v$。各因子 $\pi_v$ 在同构意义下唯一。

（2）对 $K$ 的每个连续不可约表示同构类 $\gamma$，$\pi$ 的 $\gamma$ 等型分量有限维。

（3）$\pi$ 的 $K$ 有限向量空间自然构成容许不可约 $G(A)$ 模 $\pi^K$。设 $\pi^K\simeq\bigotimes\pi_v^K$ 是定理 3 给出的 $\pi^K$ 的分解。则 $\pi_v^K$ 作为容许 $G(F_v)$ 模，同构于 $\pi_v$ 的 $K_v$ 有限向量空间。
\end{originaltheorem}

\begin{thebibliography}{9}
\bibitem[1]{c10-1} I. N. Bernshtein, \emph{All reductive p-adic groups are tame}, Functional Anal. Appl. 8 (1974), 91--93.
\bibitem[2]{c10-2} N. Bourbaki, \emph{Modules et anneaux semi-simples}, Hermann, Paris, 1958.
\bibitem[3]{c10-3} I. M. Gelfand, M. I. Graev and I. I. Pyatetskii-Shapiro, \emph{Representation theory and automorphic functions}, W. B. Saunders Co., Philadelphia, 1969.
\bibitem[4]{c10-4} H. Jacquet and R. P. Langlands, \emph{Automorphic forms on GL(2)}, Lecture Notes in Math., vol. 114, Springer-Verlag, Berlin and New York, 1970.
\bibitem[5]{c10-5} J. Tits, \emph{Reductive groups over local fields}, 本论文集，第 1 册，第 29--69 页。
\end{thebibliography}
\par\noindent Duke 大学（Duke University）。
